\chapter{Conclusions and Practical Implications}\doublespacing
\label{Chapter6}
\lhead{Chapter VI. \emph{Conclusions and Practical Implications}}

\section{Summary}
This thesis implemented Magnavis for dual-observatory magnetic monitoring and evaluated residual-based anomaly detection under a reproducible zero-historic benchmark. The study compared statistical smoothers, a session-trained attention bidirectional LSTM, and pretrained GRU and LSTM models on short-duration, long-duration, and synthetic ground-truth layouts.

\section{Study Findings}
\begin{itemize}
    \item Prediction residuals with an EWMA threshold provide a practical detection rule when the baseline is learned rather than assumed constant.
    \item On long-duration ground truth, pretrained LSTM maintains about 85\% recall for $k=2$--$5$; pretrained GRU matches at $k=3$ but falls to about 63\% for $k \ge 4$.
    \item Closed-loop autoregressive rolling does not remove the GRU high-$k$ gap on Apr-27.
    \item Direct smoothers on raw magnitude are inadequate for sustained-offset detection at moderate $k$.
    \item Perturbation-vector direction and triangulation are useful geometric tools but do not uniquely locate a source without calibration and modelling assumptions.
\end{itemize}

\section{Research Contribution}
The main contributions are: (i) an integrated open pipeline from CSV or database to forecasts, flags, and plots; (ii) bundled pretrained OBS2 checkpoints and training scripts; (iii) a documented zero-historic benchmark with aligned point-level metrics; and (iv) evidence-based guidance on model and $k$ selection for long versus short campaigns.

\section{Practical Implications}
For operational use of Magnavis on similar data, LSTM at $k \approx 3$ is the most robust setting identified for long monitoring periods. Short campaigns may use higher $k$ with careful validation. GRU deployments should avoid high $k$ on sustained events unless retrained and re-evaluated on representative labels. Multi-sensor layouts continue to support rejection of single-instrument artefacts.

\section{Study Limitations and Scope for Future Work}
\begin{itemize}
    \item Metrics are point-level on fixed CSV exports; event-level delay and duration were not optimised.
    \item Pretrained weights reflect specific training archives and may not transfer to other sites without retraining.
    \item Direction finding was not fully calibrated against known source bearings in the field.
    \item Database credentials in prototype scripts should be moved to secure configuration for production.
\end{itemize}
Future work may add event-level recall, GRU retraining studies (absolute versus delta targets), dipole inversion between observatories, and continuous integration of the benchmark after model updates.

\section{Closing Remarks}
Magnavis demonstrates that recurrent forecasting and residual thresholds can be combined in a single reproducible framework for static observatory data. The benchmark artefacts in the project repository allow independent verification of the figures and tables reported here.
